% ============================================================
% Section 160: 对称方势阱中的束缚态
% ============================================================
\section{量子力学：对称方势阱中的束缚态}
\label{sec:finite-square-well-bound}

\begin{modelbox}[题目：对称方势阱中的束缚态]
粒子在深度为 $V_0$、宽度为 $a$ 的对称方势阱中运动，求：
\begin{enumerate}
    \item 在阱口附近刚好出现一条束缚能级（$E\approx V_0$）的条件；
    \item 束缚态能级总数。
\end{enumerate}
\end{modelbox}

\begin{tipbox}[考点快速分析]
\begin{itemize}
    \item \textbf{阱口条件}：$E\to V_0^-$ 时 $\beta\to0$，阱外波函数几乎为常数
    \item \textbf{宇称分类}：偶宇称和奇宇称分别讨论
    \item \textbf{能级总数}：与无限深势阱对比，有限深势阱多一个阈上浅束缚态
    \item \textbf{一维势阱定理}：任何一维对称吸引势至少存在一个偶宇称束缚态
\end{itemize}
\end{tipbox}

\begin{derivationbox}[标准解题过程]
\textbf{1. 基本设定}

取阱底势能为零：
\[
V(x)=\begin{cases}0, & |x|\le a/2 \\ V_0, & |x|>a/2\end{cases}
\]
束缚态能量 $0<E<V_0$。引入波数
\[
k=\frac{\sqrt{2mE}}{\hbar},\quad \beta=\frac{\sqrt{2m(V_0-E)}}{\hbar},\quad k_0=\frac{\sqrt{2mV_0}}{\hbar},
\]
满足 $k^2+\beta^2=k_0^2$。

阱内：$\psi''+k^2\psi=0$；阱外：$\psi''-\beta^2\psi=0$。

\textbf{2. 阱口条件分析}

当 $E\approx V_0^-$ 时，$\beta\to0^+$，$k\approx k_0$。阱外波函数衰减极慢，几乎为常数。

\textit{偶宇称}：$\psi(x)=C\cos kx$（阱内），$\psi(x)=Ae^{-\beta|x|}$（阱外）。

$x=a/2$ 处连续性：$C\cos(ka/2)\approx A$；
导数连续性：$-Ck\sin(ka/2)\approx0$，故 $\sin(ka/2)\approx0$，即 $ka\approx2m\pi$（$m=0,1,2,\dots$）。

\textit{奇宇称}：$\psi(x)=C\sin kx$（阱内），$\psi(x)=\pm Ae^{-\beta|x|}$（阱外）。

$x=a/2$ 处连续性：$C\sin(ka/2)\approx A$；
导数连续性：$Ck\cos(ka/2)\approx0$，故 $\cos(ka/2)\approx0$，即 $ka\approx(2m+1)\pi$（$m=0,1,2,\dots$）。

统一写为 $k_0a\approx n\pi$（$n=1,2,3,\dots$），即
\begin{equation}
\boxed{\frac{2mV_0a^2}{\hbar^2}=n^2\pi^2,\quad n=1,2,3,\ldots}
\end{equation}

$n=1,3,5,\dots$ 对应偶宇称态；$n=2,4,6,\dots$ 对应奇宇称态。偶奇交替出现。

\textbf{3. 束缚态总数}

一维对称吸引势至少存在一个偶宇称束缚态。当 $0<\frac{2mV_0a^2}{\hbar^2}<\pi^2$ 时，只有一个束缚态（基态）。

每增加一个 $n^2\pi^2$ 的阈值，就新增一个束缚态。束缚态总数为
\begin{equation}
\boxed{N=1+\left\lfloor\frac{a}{\pi\hbar}\sqrt{2mV_0}\right\rfloor.}
\end{equation}

\textbf{4. 与无限深势阱的对比}

宽度为 $a$ 的无限深势阱能级 $E_n^{(\infty)}=\frac{n^2\pi^2\hbar^2}{2ma^2}$。若要求 $E_n^{(\infty)}<V_0$，则 $n<\frac{a}{\pi\hbar}\sqrt{2mV_0}$，满足条件的能级数为 $\left\lfloor\frac{a}{\pi\hbar}\sqrt{2mV_0}\right\rfloor=N-1$。

物理原因：有限深势阱的波函数可渗透到阱外，有效宽度更大，动能更低，因此在同一 $V_0$ 上限之下能``多装''一个浅束缚态。
\end{derivationbox}

\begin{formulabox}[答案汇总]
\begin{center}
\begin{tabular}{ll}
\toprule
\textbf{结论} & \textbf{结果} \\
\midrule
阱口出现新束缚态条件 & $\displaystyle\frac{2mV_0a^2}{\hbar^2}=n^2\pi^2$（$n=1,2,\ldots$） \\
束缚态总数 & $N=1+\left\lfloor\dfrac{a}{\pi\hbar}\sqrt{2mV_0}\right\rfloor$ \\
偶奇交替 & $n$ 为奇数→偶宇称；$n$ 为偶数→奇宇称 \\
\bottomrule
\end{tabular}
\end{center}
\end{formulabox}

\begin{insightbox}[物理图像]
\begin{itemize}
    \item 一维对称势阱的束缚态总数公式与势阱``尺寸'' $k_0a$ 直接相关，每跨过 $n\pi$ 阈值就多一个态。
    \item 无论势阱多浅多窄，总至少存在一个偶宇称束缚态，这是一维量子力学的重要定理（不存在零能共振的三维情况截然不同）。
    \item 在阱口附近，波函数在阱外延展极长（$\beta\to0$），有效尺寸趋于无穷，对应``临界束缚''状态。
\end{itemize}
\end{insightbox}
